\documentclass[12pt]{article}
\usepackage[utf8]{inputenc}
\usepackage[T1]{fontenc}
\usepackage{hyperref}
\title{Totoro Language Resources for K-12 Curriculum Development}
\author{Compiled by Claude Sonnet 5 (claude-sonnet-5) for Andr\'as Kornai}
\date{}
\begin{document}
\maketitle
\section{Introduction}
Totoro (Totor\'o; ISO 639-3 ttk) is a Barbacoan language of Cauca, Colombia. Some sources call it Polindara. It is critically endangered and is listed as dormant in the Endangered Languages catalogue. Wikipedia reports about 76 speakers in 2019, which conflicts with the catalogue. Spanish has displaced the language in most settings. The language is closely related to Guambiano, which is still vigorous and is sometimes called Nam Trik or Namrrik. Materials for the two languages should not be confused.
\section{Primary Educational Resources}
The search results mention a revitalization radio series and multimedia learning tools, but gave no direct links to teaching materials.
\section{Linguistic and Reference Materials}
\begin{itemize}
\item Wikipedia, Totoro language. \url{https://en.wikipedia.org/wiki/Totoro_language}
\item Endangered Languages Project, Totoro entry. \url{https://endangeredlanguages.com/lang/3495}
\item Endangered Languages Project bibliography for Totoro. \url{https://endangeredlanguages.com/lang/3495/bibliography}
\item Omniglot page on Guambiano, a related language, listed for comparison only. \url{https://www.omniglot.com/writing/guambiano.htm}
\end{itemize}
\section{Religious and Community Resources}
None found in the search results.
\section{Digital and Audio Resources}
The search results mention a radio series and multimedia tools but gave no usable links for them.
\section{Curriculum Development Context}
A heritage curriculum would have to rely on the bibliography above and on the community, and little documentation was found. The closely related Guambiano may supply comparative material.
\end{document}